% Resumo em português
A especiação de mercúrio é uma técnica analítica utilizada para identificar e
quantificar as diferentes espécies químicas desse elemento em amostras
ambientais, tais como mercúrio elementar (Hg$^0$) e metilmercúrio. O
procedimento de preparação de amostras exige o controle rigoroso de uma
sequência de etapas químicas --- derivatização com tetraetilborato de sódio,
criofocalização em nitrogênio líquido a aproximadamente $-196\,\celsius$,
aquecimento controlado do Tubo U até $230\,\celsius$ e atomização a
$700\,\celsius$ ---, de modo que a qualidade da rampa de temperatura imposta ao
Tubo U determina diretamente a correta separação e identificação das espécies.

Este trabalho apresenta o projeto e a implementação de um sistema de automação e
supervisão para a preparação de amostras em especiação de mercúrio, concebido
para substituir um sistema legado baseado em LabVIEW. O sistema adota uma
arquitetura mestre-escravo composta por um Raspberry Pi, que executa a lógica de
alto nível em Python (máquina de estados finita, controle proporcional-integral-
derivativo --- PID --- e persistência de parâmetros), um Arduino Uno operando
como módulo de aquisição de dados em tempo real, responsável pelo acionamento de
válvulas e fornos, pela leitura de termopares via interface SPI e por um
watchdog de segurança, e uma interface homem-máquina (IHM) Web desenvolvida em
React e TypeScript, com monitoramento em tempo real via WebSocket, Diagrama de
Tempos, gráficos de tendência e modos automático e manual.

O processo foi modelado em uma máquina de estados finita com quatro fases de
acionamento ($T_0$ a $T_3$) e uma matriz de atuadores (válvulas SV1 a SV5,
bomba peristáltica e fornos). O controle de temperatura emprega uma estratégia
PID mista: abaixo de $0\,\celsius$, a variável manipulada é calculada pela razão
entre a taxa de aquecimento desejada pelo operador e a taxa do sistema,
enquanto, acima desse limiar, um controlador PID em malha fechada segue um
setpoint dinâmico de rampa até $230\,\celsius$; o Forno 2 é regulado em setpoint
fixo de $700\,\celsius$. A comunicação entre o Raspberry Pi e o Arduino é
realizada por pacotes JSON em enlace serial a 4 Hz, com handshake de
parametrização e watchdog duplo que impõe o retorno ao estado seguro em caso de
perda de comunicação.

O desenvolvimento seguiu a abordagem de desenvolvimento orientado a testes, com
validação por testes unitários, de integração, funcionais e de segurança
executados em ambiente de simulação sem hardware físico. Os resultados mostram
que a arquitetura proposta atende aos requisitos funcionais e de segurança
definidos, garantindo repetibilidade da lógica de controle, persistência
auditável dos parâmetros do método e parada segura automática diante de falhas
de comunicação. A validação em bancada, com a calibração da curva de aquecimento
do Tubo U e a verificação experimental da rampa de temperatura, é apresentada
como etapa de continuidade do trabalho.

\vspace{0.5cm}
\noindent\textbf{Palavras-chave:} automação; controle PID; especiação de
mercúrio; máquina de estados finita; interface homem-máquina; sistema embarcado.
